\chapter{Chapter05 — Small Low Degree Weight PTF}
%%% generated by the blueprint pipeline from FABL.Chapter05.SmallLowDegreeWeightPTF
\section{Overview}

Declarations of the Lean module \texttt{FABL.Chapter05.SmallLowDegreeWeightPTF}.

\section{Declarations}

\begin{definition}[All-positive sign vector]
\lean{FABL.allOneSignCube}
\label{FABL.allOneSignCube}
\leanok
For each natural number $n$, the all-positive sign vector in $\{-1,1\}^n$ assigns the
value $1$ to every coordinate.
\end{definition}

\begin{definition}[Parity with one value flipped]
\lean{FABL.onePointFlippedParity}
\statementsource{Analysis-of-Boolean-Functions-by-Ryan-ODonnell}{pdf-ffbe9c7b3306-p135-b005}
\label{FABL.onePointFlippedParity}
\leanok
\uses{FABL.BooleanFunction, FABL.allOneSignCube, FABL.parityFunction}
For $k\in\mathbb N$, define a function on the $(k+1)$-dimensional Boolean cube by taking
the parity function and negating its value at the all-$1$ input, while leaving its
values unchanged at every other input.
\end{definition}

\begin{theorem}[One-point flip differs from parity]
\lean{FABL.onePointFlippedParity_ne_fullParity}
\label{FABL.onePointFlippedParity_ne_fullParity}
\leanok
\uses{FABL.allOneSignCube, FABL.onePointFlippedParity, FABL.parityFunction}
\difficulty{3}
For every integer $k\ge 0$, the Boolean function on the $(k+1)$-dimensional sign cube
that agrees with full parity except that its value is negated at the all-one input is
not equal to the full parity function.
\end{theorem}

\begin{theorem}[One-point flip differs from negative parity]
\lean{FABL.onePointFlippedParity_ne_neg_fullParity}
\label{FABL.onePointFlippedParity_ne_neg_fullParity}
\leanok
\uses{FABL.allOneSignCube, FABL.onePointFlippedParity, FABL.parityFunction}
\difficulty{4}
For every integer $k\ge 0$, the Boolean function on the $(k+1)$-dimensional sign cube
that agrees with full parity except that its value is negated at the all-one input is
not equal to the negative of the full parity function.
\end{theorem}

\begin{theorem}[Threshold degree of flipped parity]
\lean{FABL.onePointFlippedParity_isPolynomialThreshold}
\label{FABL.onePointFlippedParity_isPolynomialThreshold}
\leanok
\uses{FABL.IsPolynomialThreshold, FABL.isPolynomialThreshold_pred_of_ne_parityFunction_univ, FABL.onePointFlippedParity, FABL.onePointFlippedParity_ne_fullParity, FABL.onePointFlippedParity_ne_neg_fullParity}
\difficulty{1}
For every natural number $k$, the Boolean function on $\{-1,1\}^{k+1}$ that equals the
parity of all $k+1$ coordinates except that its value at the all-ones input is flipped
has a polynomial threshold representation of Fourier degree at most $k$.
\end{theorem}

\begin{theorem}[Real form of the one-point flip]
\lean{FABL.onePointFlippedParity_toReal}
\label{FABL.onePointFlippedParity_toReal}
\leanok
\uses{FABL.BooleanFunction.toReal, FABL.allOneSignCube, FABL.onePointFlippedParity, FABL.parityFunction, FABL.signValue}
\difficulty{4}
For every integer $k\ge 0$, the real-valued function obtained from the one-point-flipped
full parity function by viewing its signs as real numbers equals the real-valued full
parity function minus $2$ at the all-one input, and equals the real-valued full parity
function at every other input.
\end{theorem}

\begin{theorem}[Fourier coefficient of a point mass at the all-one vector]
\lean{FABL.fourierCoeff_two_mul_allOnePointMass}
\label{FABL.fourierCoeff_two_mul_allOnePointMass}
\leanok
\uses{FABL.allOneSignCube, FABL.fourierCoeff, FABL.monomial}
\difficulty{3}
Let $k$ be a nonnegative integer and let $S$ be any subset of the $k+1$ coordinates. The
Fourier coefficient at $S$ of the function that equals $2$ at the all-one vector and $0$
elsewhere on the $(k+1)$-dimensional Boolean cube is $2^{-k}$.
\end{theorem}

\begin{theorem}[Fourier coefficients of a flipped parity]
\lean{FABL.fourierCoeff_onePointFlippedParity}
\label{FABL.fourierCoeff_onePointFlippedParity}
\leanok
\uses{FABL.BooleanFunction.toReal, FABL.fourierCoeff, FABL.fourierCoeff_monomial, FABL.fourierCoeff_sub, FABL.fourierCoeff_two_mul_allOnePointMass, FABL.onePointFlippedParity, FABL.onePointFlippedParity_toReal, FABL.parityFunction_toReal}
\difficulty{1}
For a nonnegative integer $k$, let $F$ be the Boolean function on $\\{-1,1\\}^{k+1}$
that agrees with the product of all $k+1$ coordinates except at the all-ones input,
where its value is the negative of that product. For every subset $S$ of the
coordinates, its real-valued Fourier coefficient at $S$ is $1-2^{-k}$ if $S$ is the full
coordinate set, and is $-2^{-k}$ otherwise.
\end{theorem}

\begin{theorem}[Subsets below the full degree]
\lean{FABL.lowDegreeFrequencies_eq_erase_univ}
\label{FABL.lowDegreeFrequencies_eq_erase_univ}
\leanok
\difficulty{4}
For every nonnegative integer $k$, among all subsets of a set with $k+1$ elements, the
subsets of size at most $k$ are exactly all subsets except the full set.
\end{theorem}

\begin{theorem}[Low-Degree Weight of a Flipped Parity Function]
\lean{FABL.fourierWeightAtMost_onePointFlippedParity}
\label{FABL.fourierWeightAtMost_onePointFlippedParity}
\leanok
\uses{FABL.BooleanFunction.toReal, FABL.fourierCoeff_onePointFlippedParity, FABL.fourierWeight, FABL.fourierWeightAtMost, FABL.lowDegreeFrequencies_eq_erase_univ, FABL.onePointFlippedParity}
\difficulty{4}
For each natural number $k$, consider the parity function on the $(k+1)$-dimensional
sign cube that is negated at the all-positive input and unchanged elsewhere. The sum of
the squared Fourier coefficients of this function's real-valued sign representation over
subsets of size at most $k$ equals $(2^{k+1}-1)\,((1/2)^k)^2$.
\end{theorem}

\begin{theorem}[Low-degree threshold with small low-degree weight]
\lean{FABL.exists_polynomialThreshold_fourierWeightAtMost_lt_two_mul_invPow}
\label{FABL.exists_polynomialThreshold_fourierWeightAtMost_lt_two_mul_invPow}
\leanok
\uses{FABL.BooleanFunction, FABL.BooleanFunction.toReal, FABL.IsPolynomialThreshold, FABL.fourierWeightAtMost, FABL.fourierWeightAtMost_onePointFlippedParity, FABL.onePointFlippedParity, FABL.onePointFlippedParity_isPolynomialThreshold}
\difficulty{4}
For every natural number $k$, there is a Boolean function on the sign cube of dimension
$k+1$ that has a polynomial threshold representation of Fourier degree at most $k$, and
whose total squared Fourier weight on sets of size at most $k$ is strictly less than
$2(1/2)^k$.
\end{theorem}
